\documentclass[12pt]{article}

\usepackage[T1]{fontenc}
\usepackage{lmodern}
\usepackage{amsmath,amssymb,mathtools}
\usepackage[margin=1.1in]{geometry}
\usepackage{microtype}
\usepackage{setspace}

\newcommand{\R}{\mathbb{R}}
\newcommand{\X}{\mathcal{X}}
\newcommand{\dd}{\mathop{}\!\mathrm{d}}

\setstretch{1.12}
\setlength{\parindent}{0pt}
\setlength{\parskip}{0.7em}
\setlength{\abovedisplayskip}{0.9em}
\setlength{\belowdisplayskip}{0.9em}

\title{Linearization of an Affine Switched System}
\author{}
\date{}

\begin{document}
\maketitle

\section{The affine switched system}

Consider the affine system
\begin{equation}
    \dot{x}(t)=A_{\sigma(t)}x(t)+b_{\sigma(t)},
    \qquad x(t)\in\R^n.
    \tag{1}
\end{equation}
The variable $t\in\R$ is continuous time. The switching signal
$\sigma:\R_+\to\mathcal Q$ is piecewise constant and selects one mode
from the finite set $\mathcal Q$.

Let
\[
    t_0<t_1<\cdots<t_N
\]
be the actual switching times, and let
\[
    \bar t_0<\bar t_1<\cdots<\bar t_N
\]
be their nominal values. The prescribed mode on the $i$th interval is
\[
    q_i=\bar\sigma(\bar t_{i-1}).
\]
Thus,
\begin{equation}
    \sigma(t)=q_i=\bar\sigma(\bar t_{i-1}),
    \qquad t_{i-1}\leq t<t_i,
    \qquad i=1,\ldots,N.
    \tag{2}
\end{equation}
Only the interval lengths are perturbed; the order of the modes remains
fixed.

\section{Conversion to a homogeneous system}

The term $b_{\sigma(t)}$ makes (1) affine rather than homogeneous. To
include this term in a matrix multiplication, augment the state with a
constant component:
\[
    \X(t)=
    \begin{bmatrix}
        x(t)\\[3pt]
        1
    \end{bmatrix}
    \in\R^{n+1}.
\]
Since the derivative of the last component is zero,
\[
    \dot{\X}(t)=
    \begin{bmatrix}
        \dot{x}(t)\\[3pt]
        0
    \end{bmatrix}
    =
    \begin{bmatrix}
        A_{\sigma(t)}x(t)+b_{\sigma(t)}\\[3pt]
        0
    \end{bmatrix}.
\]
This can be written as
\[
    \dot{\X}(t)=
    \begin{bmatrix}
        A_{\sigma(t)} & b_{\sigma(t)}\\
        0 & 0
    \end{bmatrix}
    \begin{bmatrix}
        x(t)\\[3pt]
        1
    \end{bmatrix}.
\]
Define the augmented matrix associated with mode $q$ by
\[
    \widetilde A_q=
    \begin{bmatrix}
        A_q & b_q\\
        0 & 0
    \end{bmatrix}.
\]
The affine system is therefore equivalent to the homogeneous system
\[
    \boxed{\dot{\X}(t)=\widetilde A_{\sigma(t)}\X(t)}.
\]

\section{Propagation through the switching intervals}

Define the actual and nominal durations of interval $i$ as
\[
    \Delta t_i=t_i-t_{i-1},
    \qquad
    \Delta\bar t_i=\bar t_i-\bar t_{i-1}.
\]
Also define
\[
    F_i=\widetilde A_{q_i}
       =\widetilde A_{\bar\sigma(\bar t_{i-1})}.
\]
The matrix $F_i$ is constant throughout the $i$th interval.

For a constant matrix $F_i$, the solution over one interval is
\[
    \X(t_i)=e^{F_i\Delta t_i}\X(t_{i-1}).
\]
For the first interval,
\[
    \X(t_1)=e^{F_1\Delta t_1}\X(t_0).
\]
For the second interval,
\begin{align*}
    \X(t_2)
    &=e^{F_2\Delta t_2}\X(t_1)\\
    &=e^{F_2\Delta t_2}e^{F_1\Delta t_1}\X(t_0).
\end{align*}
Continuing interval by interval gives
\begin{equation}
    \X(t_N)
    =e^{F_N\Delta t_N}
     e^{F_{N-1}\Delta t_{N-1}}
     \cdots
     e^{F_2\Delta t_2}
     e^{F_1\Delta t_1}\X(t_0).
    \label{eq:actual}
\end{equation}
The order of these matrices is important: the transition associated
with the first interval is the rightmost matrix.

The nominal state is propagated with the nominal interval durations:
\begin{equation}
    \bar\X(t_N)
    =e^{F_N\Delta\bar t_N}
     e^{F_{N-1}\Delta\bar t_{N-1}}
     \cdots
     e^{F_2\Delta\bar t_2}
     e^{F_1\Delta\bar t_1}\bar\X(t_0).
    \label{eq:nominal}
\end{equation}

\section{Exact error equation}

Define the augmented-state error by
\[
    e(t)=\X(t)-\bar\X(t).
\]
Let $\delta t_i$ denote the perturbation of the $i$th interval duration:
\[
    \Delta t_i=\Delta\bar t_i+\delta t_i.
\]
The initial state may also be perturbed, so that
\[
    \X(t_0)=\bar\X(t_0)+e(t_0).
\]
Substitution in \eqref{eq:actual} gives
\begin{align}
    \X(t_N)
    ={}&e^{F_N(\Delta\bar t_N+\delta t_N)}
        e^{F_{N-1}(\Delta\bar t_{N-1}+\delta t_{N-1})}
        \cdots \notag\\
      &e^{F_2(\Delta\bar t_2+\delta t_2)}
        e^{F_1(\Delta\bar t_1+\delta t_1)}
        \bigl(\bar\X(t_0)+e(t_0)\bigr).
    \label{eq:perturbed-state}
\end{align}
Subtracting \eqref{eq:nominal} from \eqref{eq:perturbed-state} produces
the exact terminal error:
\begin{align}
 e(t_N)
 ={}&e^{F_N(\Delta\bar t_N+\delta t_N)}
     e^{F_{N-1}(\Delta\bar t_{N-1}+\delta t_{N-1})}
     \cdots \notag\\
   &e^{F_2(\Delta\bar t_2+\delta t_2)}
     e^{F_1(\Delta\bar t_1+\delta t_1)}
     \bigl(\bar\X(t_0)+e(t_0)\bigr) \notag\\[3pt]
   &-e^{F_N\Delta\bar t_N}
     e^{F_{N-1}\Delta\bar t_{N-1}}
     \cdots
     e^{F_2\Delta\bar t_2}
     e^{F_1\Delta\bar t_1}\bar\X(t_0).
 \tag{3}
 \label{eq:exact-error}
\end{align}

\section{First-order expansion of each transition matrix}

For a constant matrix $A$,
\[
    \frac{\dd}{\dd s}e^{As}=Ae^{As}=e^{As}A.
\]
The first-order Taylor expansion around $s=s_0$ is consequently
\[
    e^{A(s_0+\delta s)}
    =e^{As_0}+Ae^{As_0}\delta s
     +\mathcal O(\delta s^2).
\]
This identity explains the approximation used for every switching
interval.

For clarity, define the nominal transition matrices individually:
\[
    \phi_1=e^{F_1\Delta\bar t_1},\qquad
    \phi_2=e^{F_2\Delta\bar t_2},\qquad
    \ldots,\qquad
    \phi_N=e^{F_N\Delta\bar t_N}.
\]
Then each perturbed transition matrix has the approximation
\begin{align*}
    e^{F_1(\Delta\bar t_1+\delta t_1)}
    &\simeq \phi_1+F_1\phi_1\delta t_1,\\
    e^{F_2(\Delta\bar t_2+\delta t_2)}
    &\simeq \phi_2+F_2\phi_2\delta t_2,\\
    &\hspace{1.5em}\vdots\\
    e^{F_N(\Delta\bar t_N+\delta t_N)}
    &\simeq \phi_N+F_N\phi_N\delta t_N.
\end{align*}
Substituting these expressions into the exact error equation gives
\begin{align}
 e(t_N)\simeq{}&
 (\phi_N+F_N\phi_N\delta t_N)
 (\phi_{N-1}+F_{N-1}\phi_{N-1}\delta t_{N-1})
 \cdots \notag\\
 & (\phi_2+F_2\phi_2\delta t_2)
   (\phi_1+F_1\phi_1\delta t_1)
   \bigl(\bar\X(t_0)+e(t_0)\bigr) \notag\\
 &-\phi_N\phi_{N-1}\cdots\phi_2\phi_1\bar\X(t_0).
 \label{eq:expanded-factors}
\end{align}

\section{Keeping only first-order terms}

To see how the product is expanded, first consider three intervals. Up
to first order,
\begin{align*}
& (\phi_3+F_3\phi_3\delta t_3)
  (\phi_2+F_2\phi_2\delta t_2)
  (\phi_1+F_1\phi_1\delta t_1)\\[2pt]
&\quad\simeq
  \phi_3\phi_2\phi_1
  +F_3\phi_3\phi_2\phi_1\,\delta t_3\\
&\qquad
  +\phi_3F_2\phi_2\phi_1\,\delta t_2
  +\phi_3\phi_2F_1\phi_1\,\delta t_1.
\end{align*}
Terms such as
\[
    F_3\phi_3F_2\phi_2\phi_1\,\delta t_3\delta t_2
\]
are discarded because they are second order.

For $N$ intervals, the same expansion gives the nominal product
\[
    \phi_N\phi_{N-1}\cdots\phi_2\phi_1,
\]
together with one first-order term for each duration perturbation:
\begin{align*}
&F_N\phi_N\phi_{N-1}\cdots\phi_2\phi_1\,\delta t_N,\\
&\phi_NF_{N-1}\phi_{N-1}\cdots\phi_2\phi_1\,\delta t_{N-1},\\
&\phi_N\phi_{N-1}F_{N-2}\phi_{N-2}\cdots\phi_2\phi_1\,
  \delta t_{N-2},\\
&\hspace{10em}\vdots\\
&\phi_N\phi_{N-1}\cdots\phi_3F_2\phi_2\phi_1\,\delta t_2,\\
&\phi_N\phi_{N-1}\cdots\phi_3\phi_2F_1\phi_1\,\delta t_1.
\end{align*}
The matrices must remain in this order unless additional commutativity
assumptions are imposed.

Multiplying by $\bar\X(t_0)+e(t_0)$ creates three kinds of terms:
\begin{enumerate}
    \item the nominal-state term;
    \item terms linear in the initial error $e(t_0)$;
    \item terms linear in each duration perturbation $\delta t_i$.
\end{enumerate}
The nominal-state term cancels with the final term in
\eqref{eq:expanded-factors}. Products of the form
$\delta t_i e(t_0)$ are second order and are also discarded. Therefore,
\begin{align}
 e(t_N)\simeq{}&
 \phi_N\phi_{N-1}\cdots\phi_2\phi_1e(t_0)
 \notag\\[2pt]
 &+\phi_N\phi_{N-1}\cdots\phi_3\phi_2F_1\phi_1
   \bar\X(t_0)\,\delta t_1
 \notag\\
 &+\phi_N\phi_{N-1}\cdots\phi_3F_2\phi_2\phi_1
   \bar\X(t_0)\,\delta t_2
 \notag\\
 &+\phi_N\phi_{N-1}\cdots F_3\phi_3\phi_2\phi_1
   \bar\X(t_0)\,\delta t_3
 \notag\\
 &\hspace{8em}\vdots
 \notag\\
 &+\phi_NF_{N-1}\phi_{N-1}\cdots\phi_2\phi_1
   \bar\X(t_0)\,\delta t_{N-1}
 \notag\\
 &+F_N\phi_N\phi_{N-1}\cdots\phi_2\phi_1
   \bar\X(t_0)\,\delta t_N.
 \label{eq:full-linearization}
\end{align}

\section{Matrix form of the linearized model}

Define the nominal state-transition matrix explicitly as
\[
    \Phi=\phi_N\phi_{N-1}\cdots\phi_2\phi_1.
\]
Collect the duration perturbations in the vector
\[
    \delta\mathbf t=
    \begin{bmatrix}
        \delta t_1\\
        \delta t_2\\
        \vdots\\
        \delta t_{N-1}\\
        \delta t_N
    \end{bmatrix}.
\]
The sensitivity matrix is obtained by placing the coefficient of each
$\delta t_i$ in the corresponding column:
\begin{align*}
\Gamma=\Bigl[\;&
 \phi_N\phi_{N-1}\cdots\phi_3\phi_2F_1\phi_1\bar\X(t_0),
 \\[2pt]&
 \phi_N\phi_{N-1}\cdots\phi_3F_2\phi_2\phi_1\bar\X(t_0),
 \\[2pt]&
 \phi_N\phi_{N-1}\cdots F_3\phi_3\phi_2\phi_1\bar\X(t_0),
 \\[2pt]&
 \ldots,
 \\[2pt]&
 \phi_NF_{N-1}\phi_{N-1}\cdots\phi_2\phi_1\bar\X(t_0),
 \\[2pt]&
 F_N\phi_N\phi_{N-1}\cdots\phi_2\phi_1\bar\X(t_0)
 \;\Bigr].
\end{align*}
Equation \eqref{eq:full-linearization} can now be written compactly as
\begin{equation}
    \boxed{
        e(t_N)\simeq \Phi e(t_0)+\Gamma\,\delta\mathbf t
    }.
    \label{eq:final}
\end{equation}
This is the first-order terminal-error model. It separates the effect of
the initial-state error from the effect of perturbations in the switching
interval durations.

\end{document}
